\section{Entropy metrics with controlled labels}\label{sec:entropy}

This section turns a root probability map for a metric of scalar curvature
\(-1\) into a map with a large energy coefficient in a weighted scalar
inequality.  The accompanying volume term is the logarithm of the actual
volume expansion from a density no larger than the original volume density.
We will also make the gradient absolutely small and control all label
overlaps by one graph.  The last two properties allow the energy coefficient
to be chosen larger than the combinatorial quantities used later.

Write \(\exp^{[0]}(t)=t\) and
\(\exp^{[j+1]}(t)=\exp(\exp^{[j]}(t))\).  We call a function
\(\mathcal B:[1,\infty)\to[0,\infty)\) \emph{elementary} if, for some fixed
\(H\in\N\cup\{0\}\) and \(A,B>0\),
\[
 \mathcal B(q)\leq \exp^{[H]}\bigl(B(1+q)^A\bigr)\qquad(q\geq1).
\]
When the function is prescribed in a fixed dimension, its parameters
\(H,A,B\) are allowed to depend on that dimension.

\begin{theorem}[Entropy metrics]\label{thm:entropy}
Let \((M^n,g)\) be a closed connected Riemannian manifold, where \(n\geq3\)
and \(\Scal_g=-1\), and put \(\Gamma=\pi_1(M)\).  Let
\((\Omega_0,\mu_0)\) be a standard Borel probability space with a
measure-preserving \(\Gamma\)-action; the one-point space is allowed.
Suppose \(\mathcal L_0\) is a countable set with finitely many free
\(\Gamma\)-orbits and
\[
 F^{\mathrm{in}}:\widetilde M\times\Omega_0\longrightarrow
       \ell^2(\mathcal L_0),\qquad
 F^{\mathrm{in}}_\ell\geq0,\qquad
 \sum_{\ell\in\mathcal L_0}(F^{\mathrm{in}}_\ell)^2=1
\]
is equivariant, Borel in the transverse variable, and smooth along
\(\widetilde M\).  Assume the following uniformity for this input.  In
charts lifted from a fixed finite atlas on \(M\), with fixed inner buffers,
the supports are locally finite with a uniform bound on their number,
and all leafwise derivatives of the coordinate functions have uniform
bounds.  For each orbit representative \(\ell\) there is a compact set
\(K_\ell\subset\widetilde M\) such that
\(\supp F^{\mathrm{in}}_{\gamma\ell}(\,\cdot\,,\omega)
\subset\gamma K_\ell\) for all \(\gamma\in\Gamma\) and almost every
\(\omega\).  The bounds and compact sets in these assumptions may depend
on the input.

For every prescribed elementary function \(\mathcal B\) and every finite
number \(E_{\min}\), there are a standard Borel probability
\(\Gamma\)-space \((\Omega,\mu)\) extending \((\Omega_0,\mu_0)\), a
countable set \(\mathcal L\) with finitely many free \(\Gamma\)-orbits, and
equivariant leafwise data on \(\widetilde M\times\Omega\): a metric \(h\),
a function \(\phi\), a positive density \(\rho\), and a root probability
map \(F\) with labels \(\mathcal L\).  There are also \(E>E_{\min}\), an
integer \(q\geq n+1\), and a constant \(C=C(n,\mathcal B)<\infty\) such
that, on an invariant conull set of transverse parameters,
\begin{equation}\label{eq:entropy-contract}
\begin{split}
 &\rho\leq dV_g,\qquad
 W:=\log\frac{dV_h}{\rho}\geq0,\\
 &\wscal(h,\phi)\geq
       -C+2\left(W-\frac n2\log E\right)+E e_h(F),
       \qquad e_h(F)^{1/2}\leq E^{-1/100},\\
 &E>\mathcal B(q).
\end{split}
\end{equation}
Here the extension means that there is an equivariant measure-preserving
map \(\Omega\to\Omega_0\).  All output objects are Borel transversely and
smooth along leaves.  In the same lifted charts, the coefficients of
\(h\), their derivatives of every order, and the coefficients of \(h^{-1}\)
have uniform bounds; the function \(\phi\), the density ratio
\(\rho/dV_g\), its reciprocal, and their derivatives of every order have
uniform bounds.  The coordinates of \(F\) have the same derivative and
uniform local finiteness properties as the input.  There are equivariant
compact sets containing their supports, uniformly in the transverse
parameter.  All these output bounds may depend on the input and on the
chosen finite construction.

There is, in addition, a fixed \(\Gamma\)-invariant graph \(\mathcal G\)
on \(\mathcal L\), independent of the transverse parameter, of finite
maximum degree \(\Delta\).  It contains an edge between every pair of
distinct labels that are simultaneously nonzero at any point of any leaf.
We may take
\[
 q\geq \max\{n+1,\#(\Gamma\backslash\mathcal L),\Delta+1\}
\]
and a Borel equivariant coloring
\(\chi:\mathcal L\times\Omega\to\{1,\ldots,q\}\) such that labels of
graph distance one or two have different colors.  The set \(\mathcal L\)
can be taken to consist of finitely many copies of \(\mathcal L_0\).
The constant \(C\) is independent of the manifold, metric, input family,
all the uniform bounds just specified, and \(E_{\min}\).
\end{theorem}

We prove the theorem using Proposition~\ref{prop:sweep}.  The bootstrap
will fix the density \(\rho\), which then remains fixed through all later
sweeps and conformal changes.  For
\(W=\log(dV_h/\rho)\), the balance propagated after the bootstrap is
\[
 \wscal(h,\phi)-\epsilon e_h(F)-2W+n\log\epsilon\geq-C+P,
 \qquad W\geq0,\quad P\geq0.
\]
We first raise an initially small energy coefficient to a dimensional
positive value.  Subsequent sweeps have a summable scalar loss and
retain the nonnegative surplus \(P\).  At the end of a group of sweeps,
that surplus pays for a conformal change which makes the gradient
absolutely small.  The small gradient then controls how fast the fixed
graph of label overlaps can grow in the next group of sweeps.
The inverse-metric telescope and the assembly of logarithmic shortening
also occur in \cite[Lemma~6.1 and Theorem~6.3]{OpenAIScalarInessentiality2026};
the proof below establishes the forcing and entropy estimates needed here.

\subsection{Colorings and a conformal bootstrap}

We first supply the colors required for the sweeps.  Let \(\mathcal H\)
be any fixed invariant graph of maximum degree
\(\Delta_{\mathcal H}<\infty\) on a countable label set \(\mathcal L\).
Adjoin to the transverse space independent uniform priorities in
\([0,1]^{\mathcal L}\), with the \(\Gamma\)-action permuting coordinates.
This is a standard Borel probability extension.  Direct every edge
from the larger priority to the smaller, after discarding the invariant
null set on which two priorities coincide.  This set is null because the
label set is countable and the priorities have continuous laws.  For a
fixed vertex the probability of a strictly descending path of length
\(k\) is at most \(\Delta_{\mathcal H}^k/(k+1)!\): there are at most
\(\Delta_{\mathcal H}^k\) paths, and the priorities of the \(k+1\)
distinct vertices in a given descending path have the required order
with probability \(1/(k+1)!\).  The probability of arbitrarily long
such paths is therefore zero.  Since the label set is countable, on
an invariant conull set every vertex has finite descending depth.
Greedily assign to a vertex the least color in
\(\{1,\ldots,\Delta_{\mathcal H}+1\}\) unused by its lower-priority
neighbors, recursively in that finite depth.  This defines a proper
coloring.  It is Borel, since each vertex's color is determined in
some finite descending neighborhood, and it is equivariant, since
the rule only uses the invariant graph and priorities.  Applying this
to the graph square separates vertices of distance at most two.

For the initial map and throughout the bootstrap, graphs of the
required finite degree exist with input-dependent bounds.  To see
this, choose a compact set in \(\widetilde M\) whose translates cover
the universal cover.  Each label support lies in one of finitely many
fixed compact sets translated by its group label.  A support
enlargement by a bounded distance in any one of the current uniformly
controlled metrics is contained in an input-dependent compact
enlargement in the lifted metric \(g\).  Properness of the deck action
implies that only finitely many translates of any of these compact
sets can meet another.  Joining all such possible pairs, over all
transverse parameters, gives one invariant finite-degree graph for
that step.  The color construction just proved supplies the colors
of Proposition~\ref{prop:sweep}.  The graph and color choices are
made before the analytic errors for the corresponding sweep are
chosen.  This proves availability at each step of the finite
bootstrap, with input-dependent degree.  The phase argument below
will control the degree relative to the eventual energy coefficient.

Here and below all the notation of a single sweep is that of
Proposition~\ref{prop:sweep}.  In particular, its input is
\((h_0,\phi_0,F)\), its output is \((h_+,\phi_+,Y)\), and
\[
\begin{gathered}
 N^2=l\epsilon,\qquad a=\tfrac12+\theta,\qquad
 \lambda=\tfrac12-2\theta,\\
 r=\sum_i f_i^2>0,\qquad Y=f/\sqrt r,\qquad
 W_p=\sum_i\log L_i,\qquad m=\sum_i A_i^2.
\end{gathered}
\]
The weights \(\nu_i=f_i^2/r\) sum to one.  We recall the probability
statistics and their centered combinations from
\eqref{eq:sweep-statistics}:
\[
\begin{gathered}
 u_i=\frac{A_i^2}{\epsilon\nu_i},\qquad
 \sigma_i=\frac{F_i^2}{\nu_i},\qquad
 \bar u=\frac m\epsilon,\\
 M_0=a(1-r)-\frac{\bar u}{2},\qquad
 V_i=a(\sigma_i-1)+\lambda(u_i-\bar u),\qquad
 t_i=\frac{A_i^2}{f_i}.
\end{gathered}
\]
Indices of vanishing bumps are omitted.  In particular
\(\E_\nu u_i=\bar u\) and \(\E_\nu\sigma_i=1\); the ratios \(\sigma_i\)
need not be at most one.  We write \(\E\) and \(\Var\) for
this finite weighted mean and variance in the pointwise calculations.
The graph operations give
\begin{equation}\label{eq:entropy-sweep-density}
 \frac{dV_{h_+}}{dV_{h_0}}=e^{W_p},\qquad
 \phi_+=\phi_0-W_p,\qquad
 N^2 e_{h_+}(Y)\leq \frac mr.
\end{equation}
The added graph tensors are nonnegative.  Thus, as quadratic forms,
\(h_0\leq h_{i-1}\leq h_i\leq h_+\); the inverse forms reverse these
inequalities.  In particular, for every covector \(\alpha\),
\[
 |\alpha|_{h_+}\leq|\alpha|_{h_i}
       \leq|\alpha|_{h_{i-1}}\leq|\alpha|_{h_0}.
\]
The positive reserve \(\Pi_p\) is defined in
\eqref{eq:sweep-positive-reserve}, with coefficient exactly one on
\(T_0\).  In a sweep without boost its final term is omitted; in a
boosted sweep its constants \(c_*,C_*\) may depend on the fixed lower
threshold for \(\epsilon\).  The sweep inequality reads
\begin{equation}\label{eq:entropy-sweep-form}
\begin{split}
 Q&:=\frac{rN^2}{\epsilon}
 \bigl[\wscal(h_+,\phi_+)-\wscal(h_0,\phi_0)
        +\epsilon e_{h_0}(F)-\epsilon e_{h_+}(Y)\bigr]\\
 &\geq \Pi_p-4\theta^2m-\frac r\epsilon\delta_0 .
\end{split}
\end{equation}
The positive errors \(\delta_0,\eta_0\) can be prescribed independently
before each sweep.  The construction realizes those errors by choosing
the analytic parameters for each color after the geometry of the
preceding colors has been fixed.

\begin{lemma}[Conformal bootstrap]\label{lem:entropy-bootstrap}
There are constants \(\epsilon_*>0\) and \(C_{\mathrm{boot}}<4\),
depending only on \(n\), with the following property.  Starting with
the input of Theorem~\ref{thm:entropy}, finitely many sweeps without
boost, followed individually by conformal changes, produce a metric
\(h_*\), a root probability map \(F_*\), and
\(\epsilon_0\in[\epsilon_*,2\epsilon_*]\) such that
\begin{equation}\label{eq:entropy-bootstrap-output}
 \rho:=dV_{h_*}\leq dV_g,\qquad
 \Scal_{h_*}\geq-C_{\mathrm{boot}}+\epsilon_0e_{h_*}(F_*).
\end{equation}
The metric and map have the uniformity and support properties in
Theorem~\ref{thm:entropy}.  The number of sweeps may depend on the input.
\end{lemma}

\begin{proof}
We use sweeps with \(\phi_0=0\), no boost, and \(l=\theta\), where \(l>0\)
will be a sufficiently small fixed dimensional number.  Then
\(\phi_+=-W_p\).  After a sweep set
\[
 k=e^{-2W_p/(n-1)}h_+ .
\]
The scalar conformal formula and the definition of \(\wscal\) give
\begin{equation}\label{eq:entropy-bootstrap-reset}
 \Scal_k=e^{2W_p/(n-1)}
 \left(\wscal(h_+,\phi_+)+\frac{|dW_p|_{h_+}^2}{n-1}\right).
\end{equation}
The same reset contracts the volume density:
\begin{equation}\label{eq:entropy-bootstrap-density}
 dV_k=e^{-nW_p/(n-1)}dV_{h_+}
      =e^{-W_p/(n-1)}dV_{h_0}\leq dV_{h_0}.
\end{equation}
We claim that there are dimensional \(C_b\) and \(\epsilon_b>0\) such
that, whenever \(2\leq C_0\leq4\), \(0<\epsilon\leq\epsilon_b\), and
\(\Scal_{h_0}\geq-C_0+\epsilon e_{h_0}(F)\), one can choose the errors
in the sweep so that
\begin{equation}\label{eq:entropy-bootstrap-step}
 \Scal_k\geq-(C_0+C_b l\epsilon)+(1+l)\epsilon e_k(Y).
\end{equation}

We first derive the budget needed for this step.  The reset gives
\(e_k(Y)=e^{2W_p/(n-1)}e_{h_+}(Y)\).  Substituting the definition of
\(Q\) and the input scalar inequality into
\eqref{eq:entropy-bootstrap-reset}, and dropping its nonnegative
gradient term, shows that it suffices to have
\[
 \frac{\epsilon Q}{rN^2}\geq
 l\epsilon e_{h_+}(Y)+C_0(1-e^{-2W_p/(n-1)})
       -C_b l\epsilon e^{-2W_p/(n-1)}.
\]
Multiply by \(rN^2\), and use
\(N^2 e_{h_+}(Y)\leq m/r\) and \(N^2=l\epsilon\).  A sufficient budget
for \eqref{eq:entropy-bootstrap-step} is therefore
\begin{equation}\label{eq:entropy-bootstrap-budget}
 \epsilon Q\geq
 l\epsilon m+C_0l\epsilon r(1-e^{-2W_p/(n-1)})
       -C_b l^2\epsilon^2r e^{-2W_p/(n-1)}.
\end{equation}

To pay this budget when \(\epsilon\) is small, we need control of both
\(m\) and the deviation of \(r\) from one.  The sweep inequality implies
\begin{equation}\label{eq:entropy-bootstrap-coercivity}
 \epsilon Q\geq
 c_n\bigl(\epsilon^2(1-r)^2+m^2\bigr)
       -C_n\epsilon l^2m-r\delta'_0 ,
\end{equation}
where \(\delta'_0>0\) can be chosen arbitrarily small.  Here is a proof.
Fix a sufficiently small dimensional \(m_b>0\).  If \(m\geq m_b\) and
\(\epsilon_b\) is small enough, then \(\bar u=m/\epsilon\geq4a\).  Since
\[
 -M_0=ar+\frac{\bar u}{2}-a,
\]
the square \(\epsilon^2M_0^2\) controls
\(\epsilon^2(1-r)^2+m^2\).  This is one of the positive terms in
\(\epsilon\Pi_p\).

Suppose next that \(m<m_b\).  The prescription term \(D_i^0\) of
Proposition~\ref{prop:sweep} satisfies
\[
 ND_i^0+t_i=\frac{\epsilon f_i}{r}
 \left[V_i+M_0+(\lambda+\tfrac12)\bar u+
                    (\tfrac32-\lambda)u_i\right].
\]
Both displayed coefficients of \(\bar u\) and \(u_i\) are positive and
bounded below.  Moving \(V_i+M_0\) to the other side, taking positive
parts, and summing squares therefore yields
\begin{equation}\label{eq:entropy-bootstrap-t}
 \sum_i t_i^2\leq
 \frac{C\epsilon^2}{r}\bigl(\E V_i^2+M_0^2\bigr)
       +C\sum_i(ND_i^0+t_i)_+^2 .
\end{equation}
The prescription error relation in Proposition~\ref{prop:sweep} is
\(ND_i=ND_i^0-\beta_i f_i+\mathrm{err}_i\), with
\(\sum_i|\mathrm{err}_i|^2\) arbitrarily small.  Moreover, the
longitudinal component of \(\II_i\) is \(-X_i/N\), so
\[
 ND_i+t_i=-(X_i-t_i)
       +N\bigl(\tr\II_i^{\mathrm{rem}}-U_i(\phi_{i-1})\bigr).
\]
The terms \((X_i-t_i)^2\), \(N^2|\II_i^{\mathrm{rem}}|^2\), and
\(\beta_i^2f_i^2\) occur in \(\mathcal X\).
For the remaining term, choose \(m_b<1/4\).  Then
\(L_i\leq C\) and \(|d\log L_i|_{h_i}\leq C A_i|\II_i|\).
Because the metrics increase during the sweep and
\(\phi_{i-1}=-\sum_{j<i}\log L_j\), Cauchy--Schwarz gives
\[
\begin{split}
 \sum_iN^2U_i(\phi_{i-1})^2
 &\leq C m^2\sum_jN^2|\II_j|^2\\
 &\leq C' m^2\left(\mathcal X+\sum_jt_j^2\right).
\end{split}
\]
For example, the first line follows by bounding
\(|d\phi_{i-1}|_{h_{i-1}}\) by
\(C\sum_{j<i}A_j|\II_j|\), squaring with
\(\sum_jA_j^2=m\), and then summing the factor \(A_i^2\).
Choosing \(m_b\) small absorbs the last sum of \(t_j^2\) in
\eqref{eq:entropy-bootstrap-t}.  Since
\[
 m^2=\left(\sum_i f_i t_i\right)^2\leq r\sum_i t_i^2,
 \qquad
 \epsilon^2(1-r)^2\leq C\bigl(\epsilon^2M_0^2+m^2\bigr),
\]
fractions of the \(M_0^2,\E V_i^2,\mathcal X\) terms in
\(\epsilon\Pi_p\) control both required squares.  The arbitrarily
small errors in the prescription contribute \(r\delta'_0\).
The term \(-4\theta^2m\) in \eqref{eq:entropy-sweep-form} is
\(-4l^2m\).  This proves \eqref{eq:entropy-bootstrap-coercivity}.

We now compare the right side of
\eqref{eq:entropy-bootstrap-coercivity} with the required budget.
When \(m\geq m_b\), the first two terms on the right of
\eqref{eq:entropy-bootstrap-budget} are at most
\(l\epsilon(m+4r)\).  The squares
\(m^2+\epsilon^2(1-r)^2\) dominate this expression with a fixed
fraction to spare after \(l\) is chosen small and then \(\epsilon_b\)
is decreased if needed.  One may check this separately for \(r\leq2\),
where \(m\geq m_b\) controls the constant part, and for \(r>2\), where
Young's inequality controls \(\epsilon r\) by
\(\epsilon^2(r-1)^2\) and a sufficiently small fixed part of \(m^2\).
The term \(C_n\epsilon l^2m\) is covered at the same time.

When \(m<m_b\), the identity
\(L_i=(1-A_i^2)^{-1/2}\) gives \(W_p\leq C m\).
For \(r\notin(1/2,2)\), Young's inequality applied to
\(l\epsilon m(1+r)\) again gives the required domination by the two
squares, with \(l\) small.  For \(1/2<r<2\), the positive terms on the
right of \eqref{eq:entropy-bootstrap-budget}, together with the term
\(C_n\epsilon l^2m\) from the coercivity estimate, are at most
\(C_nl\epsilon m\).  Here
\[
 C_nl\epsilon m\leq \frac{c_n}{2}m^2+C'_n l^2\epsilon^2.
\]
On this region \(r e^{-2W_p/(n-1)}\) has a dimensional positive lower
bound.  We may therefore choose \(C_b\) dimensionally to cover the
last term.  These comparisons can be made strict by increasing \(C_b\)
and decreasing \(l,\epsilon_b\).  They leave a positive multiple of
\(l^2\epsilon^2r\) in the middle region, and positive square margins
outside it; dividing the latter margins by \(r\) gives a positive
lower bound for fixed \(l,\epsilon\).  We can consequently choose
\(\delta'_0\) small enough for the sweep.  This proves
\eqref{eq:entropy-bootstrap-step}.

The input map has a finite uniform energy bound in \(g\).  Thus for
some positive, possibly input-dependent \(\epsilon_{\mathrm{in}}\)
small enough,
\(\Scal_g=-1\geq-2+\epsilon_{\mathrm{in}}e_g(F^{\mathrm{in}})\).
Choose a dimensional \(\epsilon_*>0\) below \(\epsilon_b\), small
enough that \(2+2C_b\epsilon_*<4\), and decrease
\(\epsilon_{\mathrm{in}}\) further so that
\(0<\epsilon_{\mathrm{in}}<\epsilon_*\).  Iterate
\eqref{eq:entropy-bootstrap-step} until the coefficient first reaches
\(\epsilon_*\).  The output coefficient is at most
\((1+l)\epsilon_*\leq2\epsilon_*\), and
\[
 \sum_{\text{bootstrap sweeps}}l\epsilon
 =\epsilon_0-\epsilon_{\mathrm{in}}\leq2\epsilon_*.
\]
This bounds the total scalar cost by \(2C_b\epsilon_*\), independently
of the number of sweeps.  By \eqref{eq:entropy-bootstrap-density},
every reset contracts the volume density.  The final density
\(\rho=dV_{h_*}\) therefore satisfies
\eqref{eq:entropy-bootstrap-output}.  The sweep construction and
finitely many smooth conformal changes preserve all the stated
input-uniform properties.
\end{proof}

\subsection{The entropy increment and its surplus}

Fix \(\epsilon_*\) from Lemma~\ref{lem:entropy-bootstrap}.  We now use
the boost in Proposition~\ref{prop:sweep}, choosing its constant
\(\kappa>0\) sufficiently small depending on this fixed threshold.
All constants arising from the boost are consequently dimensional.
The increment estimate has a coarse region, where a fixed fraction of
the reserve remains after paying the budget, and a region near
\(m=0,r=1\), where the exact quadratic coefficient will matter.

\begin{lemma}[Entropy increment]\label{lem:entropy-increment}
There are dimensional constants \(m_0>0\), \(D_*\geq1\), \(l_0>0\),
\(\gamma_n>0\), and \(C_n<\infty\) with the following property.
For \(\epsilon\geq\epsilon_*\), take a boosted sweep with
\[
 0<l\leq l_0,\qquad \theta=D_*l\leq\tfrac1{16},
 \qquad N^2=l\epsilon.
\]
The errors can be chosen so that
\begin{equation}\label{eq:entropy-increment}
\begin{split}
 \wscal(h_+,\phi_+)-(1+l)\epsilon e_{h_+}(Y)
 \geq{}&\wscal(h_0,\phi_0)-\epsilon e_{h_0}(F)\\
 &+2W_p-n\log(1+l)-C_nl^2+P_p ,
\end{split}
\end{equation}
where the nonnegative, pointwise surplus is
\begin{equation}\label{eq:entropy-surplus-definition}
 P_p=\gamma_n\frac{\epsilon}{rN^2}\Pi_p\,
 \indic_{\{m\geq m_0\}\,\cup\,\{|1-r|\geq1/4\}}.
\end{equation}
In particular, on the indicated region \(P_p\) retains a fixed fraction
of every positive term in the sweep inequality.
\end{lemma}

\begin{proof}
The output energy bound in \eqref{eq:entropy-sweep-density} shows that,
in the units of \(Q\), it suffices to prove
\begin{equation}\label{eq:entropy-budget}
 Q\geq lm+2lrW_p-nrl\log(1+l)-C_nr l^3
             +\frac{rN^2}{\epsilon}P_p .
\end{equation}
The term \(lm\) pays for the additional \(l\epsilon e_{h_+}(Y)\).
We choose \(\delta_0\leq\epsilon l^3\) in
\eqref{eq:entropy-sweep-form}.  It remains to compare the other terms
there with the budget in \eqref{eq:entropy-budget}.

We use a trace identity for the amount of the new graph energy that
remains in the final metric.  Set
\[
 m_f:=\sum_i A_i^2(1-C_i)=N^2\sum_i|df_i|_{h_+}^2.
\]
As matrices,
\(\sum_iN^2df_i\otimes df_i=h_+-h_0\).  Hence
\begin{equation}\label{eq:entropy-trace}
 m_f=n-\tr(h_+^{-1}h_0)
       \leq n(1-e^{-2W_p/n}).
\end{equation}
The inequality is arithmetic--geometric mean applied to the positive
eigenvalues of \(h_+^{-1}h_0\), whose product is \(e^{-2W_p}\).
We also have
\begin{equation}\label{eq:entropy-volume-angle}
 L_i=(1-A_i^2)^{-1/2}.
\end{equation}

First choose a small dimensional \(m_0>0\), in particular
\(m_0\ll\epsilon_*\) and \(m_0<1/4\).  We claim that
\begin{equation}\label{eq:entropy-coarse-coercivity}
 T_0+\epsilon M_0^2\geq c_n(m+rW_p)
 \quad\text{if }m\geq m_0
       \text{ or if }m<m_0,\ |1-r|\geq\tfrac14.
\end{equation}
For \(m\geq2n\), the \(C_i\) terms of \(T_0\) give
\(T_0\geq m-m_f\geq m-n\geq m/2\).
For \(m_0\leq m\leq2n\), first suppose that every \(A_i^2\) is at most
a small number \(\eta=\eta(n,m_0)>0\).  Then
\[
 2W_p=\sum_i-\log(1-A_i^2)\leq(1+C\eta)m.
\]
By \eqref{eq:entropy-trace},
\[
 m-m_f\geq m-n\bigl(1-e^{-(1+C\eta)m/n}\bigr)\geq c_n m.
\]
For the last inequality, the same expression with \(\eta=0\) divided
by \(m\) has a positive minimum on \([m_0,2n]\), and then \(\eta\)
is chosen small.  If instead some \(A_i^2>\eta\), the self term of
that index in \(T_0\) is at least \(c_n(r+u_i)\).  It gives a
dimensional positive bound if \(r\) is bounded below.  If \(r\) is
small, \(M_0=a(1-r)-\bar u/2\) is bounded away from zero unless \(\bar u\)
lies in a fixed compact subinterval of \((0,\infty)\).  In the former
case \(\epsilon M_0^2\geq c_n\epsilon_*\).  In the latter,
\(\epsilon=m/\bar u\leq C_n\) and
\(u_i=A_i^2/(\epsilon\nu_i)\geq c_n\), since \(\nu_i\leq1\).
Each bound controls \(m\leq2n\).  This proves the \(m\) part for
\(m\geq m_0\).

For the volume part, indices with \(A_i\geq1/4\) satisfy
\[
 A_i^2(L_i-1)2ar\geq c\,r(L_i-1)\geq c\,r\log L_i.
\]
For all other indices, \(\log L_i\leq C A_i^2\), so their contribution
to \(rW_p\) is at most \(Crm\).  If \(r\leq2\), the bound for \(m\)
already covers this.  If \(r>2\), then
\(|M_0|\geq c(r+\bar u)\), and
\(\epsilon M_0^2\geq c\epsilon r\bar u=c\,rm\).
This proves \eqref{eq:entropy-coarse-coercivity} when \(m\geq m_0\).
When \(m<m_0\), we have \(\bar u\leq m_0/\epsilon_*\) small and
\[
 W_p=\frac m2+O\left(\sum_iA_i^4\right)\leq C m.
\]
If \(|1-r|\geq1/4\), choosing \(m_0\) still smaller gives
\(|M_0|\geq c|1-r|\).  The quantity
\(\epsilon(1-r)^2\) controls \(m+rm\): for bounded \(r\) use
\(\epsilon\geq\epsilon_*\) and \(m\leq m_0\), and for \(r\geq2\)
use \((1-r)^2\geq r^2/4\).  This proves the claim in the remaining
case.

On the region of \eqref{eq:entropy-coarse-coercivity}, the positive
reserve \(\Pi_p\) consequently controls \(m+rW_p\).  The positive terms
\(lm+2lrW_p\) in \eqref{eq:entropy-budget} use at most \(C_nl\Pi_p\),
and \(4\theta^2m\) uses at most \(C_nD_*^2l^2\Pi_p\).  For \(l\)
sufficiently small after \(D_*\), a fixed fraction of \(\Pi_p\)
remains.  Choose \(\gamma_n\) below that fraction.  The error term
\(r\delta_0/\epsilon\) is covered by \(C_nr l^3\).  This proves
\eqref{eq:entropy-budget} on the region where the surplus is retained.

It remains to prove the budget when
\begin{equation}\label{eq:entropy-small-region}
 m<m_0,\qquad |1-r|<\tfrac14.
\end{equation}
Here \(T_0\geq c_nm^2\).  Indeed, expanding the exponential in
\eqref{eq:entropy-trace} and using
\(W_p=m/2+O(\sum_iA_i^4)\) gives
\[
 \sum_i A_i^2 C_i=m-m_f
       \geq c_n m^2-C_n\sum_iA_i^4
\]
after decreasing \(m_0\).  The self terms
\(\sum_i A_i^2(L_i-1)2ar\) are at least
\(c_n\sum_iA_i^4\), by \(r>3/4\) and \(L_i-1\geq A_i^2/2\).
Adding a sufficiently small fraction of the first estimate to the
second proves \(T_0\geq c_nm^2\).
Since \(W_p\leq Cm\) on this region, this controls the budget whenever
\(m\geq Kl\), for a fixed sufficiently large dimensional \(K\).
In making this conclusion, first fix \(K\); after \(D_*\) is chosen
below, take \(l\) small enough that \(4D_*^2l^2m\) is also a small
fraction of \(m^2\).

We are left with the critical regime \(m\leq Kl\), still under
\eqref{eq:entropy-small-region}.  Here
\(W_p=m/2+O(m^2)\) and \(3/4<r<5/4\), so the budget expands as
\[
 lm+2lrW_p-nrl\log(1+l)
 =2lm-nl^2+O_n(l^3)+O_n(l^2)|r-1|.
\]
The target quadratic term is covered exactly by
\[
 \frac{m^2}{n}-2lm+nl^2
       =\left(\frac m{\sqrt n}-\sqrt n\,l\right)^2\geq0.
\]
We must therefore recover \(m^2/n\) from \(T_0\), while paying the
displayed \(r-1\) error and the errors in that recovery.  This is where
the coefficient one on \(T_0\) in \(\Pi_p\) is needed.

In particular \(e^{W_p}=1+O(m)\).  Metric addition gives
\(h_0\leq h_j\leq h_+\) as quadratic forms.
Every eigenvalue of \(h_0^{-1}h_+\) is at least one, and their product
is \(e^{2W_p}\).  Each eigenvalue is therefore at most that product,
which proves the comparison
\[
 h_0\leq h_j\leq h_+\leq e^{2W_p}h_0=(1+O(m))h_0.
\]
Normalize the directions \(e_i\) in the \(h_0\) inner product and write
the resulting unit vectors as \(\bar e_i\).  Put
\(\xi_{ij}=(\bar e_i\cdot\bar e_j)^2\).  Before approximating the metric
losses, their exact earlier- and later-color telescopes are
\[
 B_i=\sum_{j<i}(Ndf_j(e_i))^2,\qquad
 C_i=\frac1{A_i^2}\sum_{j>i}A_j^2(Ndf_i(e_j))^2
       \quad(A_i>0).
\]
The first identity follows from
\(h_{i-1}=h_0+N^2\sum_{j<i}df_j\otimes df_j\) and
\(|e_i|_{h_{i-1}}=1\).  The second follows by telescoping the rank-one
inverse formula
\[
 h_j^{-1}=h_{j-1}^{-1}-A_j^2 e_j e_j^{\mathsf T}
\]
from \(h_i\) to \(h_+\), and using
\(N^2|df_i|_{h_i}^2=A_i^2\).  Now
\(Ndf_j=A_jL_j e_j^\flat\), with the dual taken in \(h_{j-1}\).
The common metric comparison and \(L_j=1+O(m)\) show that each summand
in the exact formulas differs from its \(h_0\) expression by
\(O(m)A_j^2\).  Their sums therefore give, for \(A_i\ne0\),
\begin{equation}\label{eq:entropy-BC-expansion}
 B_i=\sum_{j<i}A_j^2\xi_{ij}+O(m^2),\qquad
 C_i=\sum_{j>i}A_j^2\xi_{ij}+O(m^2).
\end{equation}

Inserting \eqref{eq:entropy-BC-expansion} into \(T_0\) yields
\begin{equation}\label{eq:entropy-sharp-T}
 T_0\geq
 \tr\left[\left(\sum_i A_i^2\bar e_i\bar e_i^{\mathsf T}\right)^2\right]
 -C_n\left(m^3+m^2|r-1|
              +m\sum_iA_i^2(1-\sigma_i)_+\right).
\end{equation}
To see the leading coefficient, the \(B_i\) and \(C_i\) terms each
give one copy of
\(\sum_{i<j}A_i^2A_j^2\xi_{ij}\).  In the \(B_i\) terms use
\(4a^2\geq1\), \(B_i\geq0\); replacing \(\sigma_i\) by one costs at most
\(Cm\sum_iA_i^2(1-\sigma_i)_+\).  The self terms use
\(L_i-1\geq A_i^2/2\) and \(2a\geq1\).  They give
\(\sum_iA_i^4\), up to the displayed \(r-1\) and \(\sigma_i-1\) costs.
The errors in \eqref{eq:entropy-BC-expansion} sum to \(O(m^3)\).
This proves \eqref{eq:entropy-sharp-T}.  The trace in that estimate
is at least \(m^2/n\), since the matrix is positive semidefinite
with trace \(m\).

The last error in \eqref{eq:entropy-sharp-T} is paid for by the
centered squares.  From
\(a(\sigma_i-1)=V_i-\lambda(u_i-\bar u)\), Cauchy--Schwarz and Young's
inequality give, for every fixed \(c'>0\),
\begin{equation}\label{eq:entropy-v-error}
\begin{split}
 m\sum_iA_i^2(1-\sigma_i)_+
 &\leq C m\epsilon\,\E\bigl[u_i(|V_i|+|u_i-\bar u|)\bigr]\\
 &\leq c'\epsilon\bigl(\E V_i^2+\theta\Var u\bigr)
   +C_{c'}\frac{m^2}{\theta}\epsilon(\bar u^2+\Var u).
\end{split}
\end{equation}
Choose \(c'\) to use a small fraction of the squares in \(\Pi_p\).
The coefficient of \(\epsilon\Var u\) in the last term, relative to
\(\theta\epsilon\Var u\), is at most \(C_{c'}K^2/D_*^2\), because
\(m\leq Kl\) and \(\theta=D_*l\).  Choose \(D_*\) sufficiently large
after \(K\) so that this term is absorbed as well.  The remaining term
is
\[
 C_{c'}\frac{m^2}{\theta}\epsilon \bar u^2
       =C_{c'}\frac{m^4}{\theta\epsilon}=O_n(l^3).
\]
Here \(\epsilon\geq\epsilon_*\), and \(K,D_*,\epsilon_*\) are already
fixed dimensionally.  Also
\[
 |r-1|\leq C(|M_0|+\bar u).
\]
Consequently a term \(O_n(l^2)|r-1|\) is bounded by a small fraction
of \(\epsilon M_0^2\) plus \(O_n(l^3)\): use Young's inequality on
\(l^2|M_0|\), and use \(\bar u\leq Kl/\epsilon_*\) on the other part.
The errors \(m^3\), \(m^2|r-1|\), and \(4D_*^2l^2m\) in the present
region are therefore covered with a remaining error \(O_n(l^3)\).

The just established mean-square estimate also covers the
\(O_n(l^2)|r-1|\) term in the critical budget.  The recovered
\(m^2/n\) then covers \(2lm-nl^2\) by the displayed square.
Since \(r\) is bounded above and below, the remaining \(O_n(l^3)\)
terms are of the form \(C_nr l^3\).  This proves
\eqref{eq:entropy-budget} in \eqref{eq:entropy-small-region}.

The choices just made have a definite order: first \(m_0\) after
\(\epsilon_*\), then \(K\), then \(D_*\), and finally \(l_0>0\)
small enough for all the comparisons and for \(D_*l_0\leq1/16\).
The errors in each sweep are chosen afterward.  Thus every constant
in \eqref{eq:entropy-increment} is dimensional.
\end{proof}

\begin{lemma}[Bounds supplied by the surplus]\label{lem:entropy-surplus}
For every sweep of Lemma~\ref{lem:entropy-increment},
\begin{equation}\label{eq:entropy-pass-volume}
 W_p\leq C_n\epsilon^{1/4}(\epsilon+P_p)^{3/4}.
\end{equation}
There is also a dimensional \(\epsilon_{\mathrm{exit}}\) such that,
if \(\epsilon\geq\epsilon_{\mathrm{exit}}\) and \(r\leq1/4\), then
\begin{equation}\label{eq:entropy-last-surplus}
\begin{split}
 P_p&\geq c_n\frac{\epsilon}{lr}
                       \left(1+\frac{m^2}{\epsilon^2}\right),\\
 \sum_i\bigl(D_i^2+|\II_i|^2+|d\log L_i|_{h_i}^2\bigr)
       &\leq C_n\frac{P_p}{\theta}.
\end{split}
\end{equation}
The errors in the prescription may simultaneously be chosen to have
\(\sum_i|\mathrm{err}_i|^2\leq1\).
\end{lemma}

\begin{proof}
If \(\bar u=m/\epsilon\) is larger than a fixed dimensional number, then
\(m\geq m_0\), so the surplus retains the mean square and boost.
Since \(|M_0|\geq c(r+\bar u)\) in this range and \(l\leq1\),
\[
 P_p\geq c_n\epsilon\,
 \frac{(r+\bar u)^2+c_*r^6\indic_{\{r\geq C_*\}}}{r}
       \geq c'_n\epsilon \bar u^{5/3}.
\]
For the last inequality, if \(r\leq \bar u^{1/3}\) use \(\bar u^2/r\);
if \(r\geq \bar u^{1/3}\) use \(r^5\), taking the threshold for \(\bar u\)
large enough that \(\bar u^{1/3}\geq C_*\).  It follows in all cases that
\begin{equation}\label{eq:entropy-m-surplus}
 m\leq C_n\epsilon(1+P_p/\epsilon)^{3/5}.
\end{equation}
If any \(A_i>1/2\), then \(m>1/4>m_0\), and the self terms in \(T_0\)
give
\[
 P_p\geq c_n\sum_{A_i>1/2}(L_i-1).
\]
Indeed their contribution to \(T_0\) is at least
\(c r\sum_{A_i>1/2}(L_i-1)\), and
\(\epsilon/(rN^2)=1/(lr)\geq1/r\).
For \(A_i\leq1/2\), \(\log L_i\leq C A_i^2\).
There are at most \(4m\) indices with \(A_i>1/2\).  Concavity of
\(\log(1+s)\) on those indices, together with the last estimate,
therefore gives
\begin{equation}\label{eq:entropy-log-surplus}
 W_p\leq C_n m\bigl(1+\log(1+P_p/m)\bigr)
\end{equation}
when \(m>0\); when \(m=0\), \(W_p=0\).
To justify the bound when the number of those indices is smaller
than \(4m\), use that \(s\log(1+b/s)\) is increasing in \(s>0\).
Put \(x=m/\epsilon\), \(y=P_p/\epsilon\).  The elementary bound
\(1+\log(1+t)\leq C(1+t)^{1/4}\) for \(t\geq0\), applied with
\(t=y/x\), makes the right side of
\eqref{eq:entropy-log-surplus}, divided by \(\epsilon\), at most
\[
 C_n x^{3/4}(x+y)^{1/4}
 \leq C_n(1+y)^{7/10}
 \leq C_n(1+y)^{3/4},
\]
where \eqref{eq:entropy-m-surplus} was used in the middle inequality.
This proves \eqref{eq:entropy-pass-volume}.

Now suppose \(r\leq1/4\).  The surplus retains all positive terms
because \(|1-r|\geq3/4\).  Outside a fixed compact subinterval of
\(\bar u\in(0,\infty)\), the expression \(M_0=a(1-r)-\bar u/2\) satisfies
\(M_0^2\geq c(1+\bar u^2)\).  On that compact interval, for \(\epsilon\)
sufficiently large,
\[
 T_0\geq m-m_f\geq\epsilon \bar u-n
       \geq c_n\epsilon(1+\bar u^2)
\]
by \eqref{eq:entropy-trace}.  Multiplying the appropriate retained
term by \(\gamma_n/(lr)\) proves the first inequality in
\eqref{eq:entropy-last-surplus}.

For the second, the retained derivative terms imply
\begin{equation}\label{eq:entropy-X-retained}
 \mathcal X\leq C_n N^2P_p,\qquad
 \frac{\epsilon^2}{rN^2}
       \bigl(\E V_i^2+\theta\Var u\bigr)\leq C_nP_p.
\end{equation}
Also
\[
 \sum_i t_i^2=\frac{\epsilon^2}{r}(\bar u^2+\Var u).
\]
The first inequality of \eqref{eq:entropy-last-surplus} controls
the \(\bar u^2\) part of this sum divided by \(N^2\), and the variance
term in \eqref{eq:entropy-X-retained} controls the other part with
a factor \(1/\theta\).  Thus
\(\sum_i t_i^2/N^2\leq C_nP_p/\theta\).  Since \(\mathcal X\) contains
\((X_i-t_i)^2+N^2|\II_i^{\mathrm{rem}}|^2+N^2|d\log L_i|^2\), this proves the
claimed bounds for \(\II_i\) and \(d\log L_i\).

For \(D_i\), use
\[
 ND_i^0=\frac{\epsilon f_i}{r}
             \left(a(\sigma_i-r)+\frac{u_i}{2}\right),\qquad
 \E \sigma_i^2\leq C\bigl(1+\E V_i^2+\Var u\bigr).
\]
Here the second inequality follows by solving the definition of \(V_i\)
for \(\sigma_i-1\).  Since \(r\leq1/4\),
\[
 \sum_i(D_i^0)^2\leq
 \frac{C\epsilon^2}{rN^2}
       \bigl(1+\bar u^2+\E V_i^2+\Var u\bigr)
 \leq C_nP_p/\theta.
\]
The relation \(ND_i=ND_i^0-\beta_i f_i+\mathrm{err}_i\) finishes the
estimate: \(\sum_i\beta_i^2f_i^2\leq\mathcal X\), and we choose
\(\sum_i|\mathrm{err}_i|^2\leq1\).  Its contribution \(1/N^2\) is
bounded by \(C_nP_p\) using the first inequality of
\eqref{eq:entropy-last-surplus}, \(\epsilon\geq1\), and \(r\leq1/4\).
\end{proof}

\subsection{Accumulation and an absolute gradient bound}

The step size has two jobs: the scalar losses \(O_n(l^2)\) must be
summable, and the number of support enlargements before reaching \(E\)
must grow at most cubically in \(\log E\).  Since one sweep changes
\(\log(1+\epsilon)\) by a quantity comparable to \(l\), we take
\begin{equation}\label{eq:entropy-step-schedule}
 l(\epsilon)=\frac{l_*}{(1+\log(1+\epsilon))^2},
\end{equation}
where the dimensional \(l_*>0\) is chosen below \(l_0\) from
Lemma~\ref{lem:entropy-increment}.  Thus all the parameter choices in
that lemma, including \(\theta=D_*l\), apply.

The fixed exponents in the next two lemmas leave room, respectively,
for the number of sweeps, the conformal scalar cost, the outgoing
volume estimate, and the final gradient target:
\[
 0.26=\tfrac14+0.01,\qquad
 0.26+0.04<\frac{1+0.32}{4},\qquad
 0.26+\tfrac34(0.04)<0.32,\qquad
 0.02>\tfrac1{100}.
\]
The estimates below show how each margin is used.

\begin{lemma}[Accumulation of entropy passes]\label{lem:entropy-accumulation}
Suppose at parameter \(\epsilon_{\mathrm{in}}\geq\epsilon_*\) that
\[
 \wscal(h,\phi)\geq-C+2W-n\log\epsilon_{\mathrm{in}}
                  +\epsilon_{\mathrm{in}}e_h(F)+P_{\mathrm{in}},
 \qquad W=\log(dV_h/\rho)\geq0,\quad P_{\mathrm{in}}\geq0.
\]
Successive sweeps of Lemma~\ref{lem:entropy-increment}, with
\(\epsilon_{\mathrm{new}}=(1+l(\epsilon))\epsilon\), preserve this
form of inequality: at the final parameter \(E\) the surplus is
\(P=P_{\mathrm{in}}+\sum_pP_p\), the volume logarithm is
\(W_{\mathrm{new}}=W+\sum_pW_p\), and the increase of the scalar
constant is \(C_n\sum_p l(\epsilon_p)^2\).  Along the entire increasing
sequence of parameters,
\begin{equation}\label{eq:entropy-summability}
 \sum_p l(\epsilon_p)^2\leq C_n,\qquad
 \#\{p:\epsilon_p<E\}\leq C_n(1+\log^3(1+E)).
\end{equation}
If the initial volume logarithm is zero, then, for all sufficiently
large dimensional \(E\),
\begin{equation}\label{eq:entropy-phase-volume}
 W_{\mathrm{new}}\leq E^{0.26}(E+P)^{0.75}.
\end{equation}
The same assertion holds if the initial data have parameter \(E_0\),
surplus \(P_0\), and
\[
 W\leq E_0^{0.32}(E_0+P_0)^{0.75},
\]
provided \(E\geq E_0\) is sufficiently large dimensionally and
\(E_0^{0.32}\leq\tfrac12E^{0.26}\).
\end{lemma}

\begin{proof}
The accumulation assertion follows by adding
\eqref{eq:entropy-increment}.  The energy coefficients telescope
multiplicatively, and
\(\sum_p\log(1+l(\epsilon_p))=\log(E/\epsilon_{\mathrm{in}})\).
The density identity in \eqref{eq:entropy-sweep-density} gives the
sum of the \(W_p\).

Put \(t_p=\log(1+\epsilon_p)\).  Since
\(\epsilon_p\geq\epsilon_*>0\) and \(l_*\) is fixed small,
\[
 t_{p+1}-t_p
 =\log\left(1+l(\epsilon_p)\frac{\epsilon_p}{1+\epsilon_p}\right)
 \asymp_n l(\epsilon_p)
 =\frac{l_*}{(1+t_p)^2}.
\]
Comparison with the integrals of \((1+t)^2\) and \((1+t)^{-2}\)
respectively proves the number bound and
\[
 \sum_p l(\epsilon_p)^2
 \leq C_n l_*\int_{t_0}^{\infty}\frac{dt}{(1+t)^2}\leq C_n.
\]
The comparison is valid on each interval \([t_p,t_{p+1}]\), since
its endpoints differ by a uniformly bounded small number.

Let \(k\) be the number of sweeps in the group under consideration.
By Lemma~\ref{lem:entropy-surplus}, since
\(\epsilon_p\leq E\) and \(P_p\leq P\),
\[
 \sum_pW_p\leq C_n k E^{1/4}(E+P)^{3/4}.
\]
For all sufficiently large \(E\),
\(C_n k\leq \tfrac12 E^{0.01}\) by
\eqref{eq:entropy-summability}.  This proves
\eqref{eq:entropy-phase-volume} for zero initial volume logarithm.
In the other case \(P\geq P_0\) and \(E\geq E_0\), so the initial
term is at most \(E_0^{0.32}(E+P)^{0.75}\).  The stated condition on
\(E_0\) pays for this with the other half of
\(E^{0.26}(E+P)^{0.75}\).
\end{proof}

Call a finite group of these sweeps a \emph{phase}.  There is no
conformal change inside a phase.  Each sweep adds the nonnegative
tensor in \eqref{eq:sweep-output}; hence, if \(h^0\) is its starting
metric and \(h_+\) its final metric, the successive metrics and their
inverses satisfy
\[
 h^0\leq h^{(1)}\leq\cdots\leq h_+,\qquad
 (h_+)^{-1}\leq\cdots\leq(h^{(1)})^{-1}\leq(h^0)^{-1}.
\]
In particular, the starting metric is an admissible reference metric
for every sweep in the phase, and
\(e_{h_+}(Y^0)\leq e_{h^0}(Y^0)\) for its starting map \(Y^0\).

We now finish a phase with an absolute gradient bound.  Normalizing
the last bumps alone divides by \(\sqrt r\), which can make the
gradient large where \(r\) is small.  We first keep a small copy of
the starting map in the normalization; its gradient is controlled by
the preceding metric comparison.  A conformal enlargement confined
to the small-\(r\) region then reduces the remaining gradient there,
and the last sweep's surplus pays its cost.

\begin{lemma}[Exit from a phase]\label{lem:entropy-phase-exit}
Consider a phase starting with root map \(Y^0\) and metric \(h^0\), and
let \(K_0\) be a finite uniform bound for \(e_{h^0}(Y^0)^{1/2}\).
Suppose its last sweep has input parameter \(\epsilon\), output
parameter \(E=(1+l(\epsilon))\epsilon\), and output
\((h_+,\phi_+,Y)\).  Suppose, with \(C\geq0\) and \(P\geq P_p\) including the
surplus from that last sweep, that
\begin{equation}\label{eq:entropy-before-exit}
\begin{split}
 \wscal(h_+,\phi_+)&\geq-C+2W-n\log E+E e_{h_+}(Y)+P,\\
 W&=\log(dV_{h_+}/\rho)\geq0,\qquad
 W\leq E^{0.26}(E+P)^{0.75}.
\end{split}
\end{equation}
If \(E\) is sufficiently large in terms of \(n,K_0\), there are a
metric \(\widehat h\), a weight \(\widehat\phi\), and a root map \(Z\)
on two copies of the labels such that
\begin{equation}\label{eq:entropy-after-exit}
\begin{split}
 \wscal(\widehat h,\widehat\phi)
 &\geq-(C+C_n)+2\widehat W-n\log E
                      +E e_{\widehat h}(Z)+P',\\
 \widehat W&=\log(dV_{\widehat h}/\rho)\geq0,\qquad
 \widehat W\leq E^{0.32}(E+P')^{0.75},\\
 e_{\widehat h}(Z)^{1/2}&\leq E^{-1/100}.
\end{split}
\end{equation}
Here \(P'\geq0\), and the added scalar cost \(C_n\) is independent
of \(K_0\) and of all other instance constants.  The changes preserve
the uniform smoothness and support properties of
Theorem~\ref{thm:entropy}.
\end{lemma}

\begin{proof}
In this proof all derivatives, unless otherwise indicated, use \(h_+\).
Write \(r=\sum_i f_i^2\) for the last sweep, set \(s=\sqrt r\), and put
\[
 c_E=\frac1{\sqrt E(1+K_0)},\qquad
 Z=\frac{(f,c_EY^0)}{\sqrt{r+c_E^2}}.
\]
The two vectors in the numerator use disjoint copies of the labels.
The map is a smooth root probability map because its denominator is
at least \(c_E>0\).  Writing the copies as \((i,0)\) and \((i,1)\),
we call \(i\) their common parent label at the start of the phase;
the sweeps within the phase do not change the label set.  The first
copy is supported on the last bumps and the second on the starting
map.  This parent assignment will be used to compare support overlaps
between phases.  By \eqref{eq:entropy-trace},
\[
 N^2\sum_i|df_i|_{h_+}^2=m_f\leq n,\qquad
 |ds|_{h_+}^2\leq n/N^2.
\]
Monotonicity of the phase metrics also gives \(e_{h_+}(Y^0)\leq K_0^2\).

Since \(Y\) and \(Y^0\) are unit vectors in orthogonal spaces, their
derivatives are perpendicular to their respective vectors.  Differentiating
\(Z=(sY,c_EY^0)/\sqrt{s^2+c_E^2}\) thus gives the exact identity
\[
 e(Z)=\frac r{r+c_E^2}e(Y)+\frac{c_E^2}{r+c_E^2}e(Y^0)
           +\frac{c_E^2}{(r+c_E^2)^2}|ds|^2.
\]
For \(r\geq1/4\), it follows that
\[
 E(e(Z)-e(Y))\leq
 \frac{Ec_E^2K_0^2}{r+c_E^2}
       +\frac{Ec_E^2n}{N^2(r+c_E^2)^2}
 \leq 4+\frac{16n}{N^2}\leq C_n
\]
when \(E\) is large, since \(N^2=l\epsilon\) tends to infinity under
\eqref{eq:entropy-step-schedule}.  For \(r\leq1/4\), the norm of the
derivative of normalization is at most \(1/\sqrt{r+c_E^2}\), so
\begin{equation}\label{eq:entropy-Z-energy}
 e(Z)\leq\frac{n/N^2+c_E^2K_0^2}{r+c_E^2}.
\end{equation}
The first inequality of \eqref{eq:entropy-last-surplus} now implies
\[
 E e(Z)\leq \frac{C_n}{l(r+c_E^2)}
       \leq \frac{C_n}{\epsilon}P_p.
\]
Thus replacing \(Y\) by \(Z\) in
\eqref{eq:entropy-before-exit} costs a dimensional constant on
\(\{r\geq1/4\}\) and an arbitrarily small fraction of \(P_p\) on
\(\{r\leq1/4\}\), once \(E\) is dimensionally large.  In particular,
no factor \(K_0\) occurs in this scalar cost.

Choose a smooth function \(u=u(r)\geq0\), supported in
\(\{r\leq E^{-0.32}\}\), with
\begin{equation}\label{eq:entropy-u-choice}
 e^u\asymp
 \min\bigl(E^{0.02},\max(1,E^{-0.16}/\sqrt r)\bigr),
 \qquad \left|\frac{du}{dr}\right|\leq\frac C r.
\end{equation}
For example, smooth, with bounded-width transitions on the logarithmic
scale, the function obtained by truncating
\(-\tfrac12\log r-0.16\log E\) to the interval
\([0,0.02\log E]\).  The smoothing may vanish whenever the untruncated
function is nonpositive, and changes its value by a bounded amount;
this gives both the support assertion and the comparison in
\eqref{eq:entropy-u-choice}.  Define
\[
 \widehat h=e^{2u}h_+,\qquad
 \widehat\phi=\phi_+-(n-1)u,\qquad
 \widehat W=W+nu.
\]
In particular \(dV_{\widehat h}=e^{nu}dV_{h_+}\), so this is the
claimed volume logarithm.  Differentiating \(r\) gives
\[
 |dr|\leq2\sqrt r\left(\sum_i|df_i|^2\right)^{1/2}
       \leq\frac{2\sqrt{nr}}N,\qquad
 |du|\leq\frac{C_n}{N\sqrt r}.
\]
The coefficient \(n-1\) in the new weight cancels the Laplacian of
\(u\) in the weighted conformal formula.  With
\(b=\nabla_{h_+}\phi_+\), that formula is
\begin{equation}\label{eq:entropy-conformal-exit}
 \wscal(\widehat h,\widehat\phi)
 =e^{-2u}\left[\wscal(h_+,\phi_+)+2b(u)-(n-1)|du|^2\right].
\end{equation}

With the fixed density \(\rho\), a graph sweep changes \(W\) by
\(+W_p\) and \(\phi\) by \(-W_p\), so it preserves \(\phi+W\).
The present exit instead satisfies
\[
 \widehat\phi+\widehat W=\phi_++W+u.
\]
Thus after an exit \(W\) need not equal \(-\phi\), and its volume
bound supplies no bound for the gradient of the weight.  We need a
bound on \(b(u)\) which uses no bound on the size of \(b\).
It is enough to estimate this on the support of \(u\), which lies in
\(\{r\leq1/4\}\) for large \(E\).  At such a point view all the graphs
of the last sweep in the common
Euclidean extension of \((T_x\widetilde M,h_0)\) by their graph
directions, one for each active color at that point.  Write
\(\mathbf n_i\) for the successive unit normals.  They are
orthonormal: a later
normal lies in the preceding graph plus its new vertical direction
and is perpendicular to every earlier normal.  In this space, each
gradient is embedded through its graph and padded by zero in the
unused later directions.  Writing \(b_i\) for the embedded gradient
of \(\phi_i\), orthogonal projection to the new graph gives the exact
recurrence
\[
 b_i=b_{i-1}-\langle b_{i-1},\mathbf n_i\rangle\mathbf n_i
             -\nabla_{h_i}\log L_i,
\]
where the last gradient is embedded in the same space.  The normal
coefficient is
\(\langle b_{i-1},\mathbf n_i\rangle=-U_i(\phi_{i-1})\), and
its magnitude satisfies
\[
 |U_i(\phi_{i-1})|\leq |D_i|+\sqrt n\,|\II_i|,
\]
because \(D_i=\tr\II_i-U_i(\phi_{i-1})\).
By orthonormality and \eqref{eq:entropy-last-surplus}, the norm of
the sum of these normal corrections is at most
\(C_n\sqrt{P_p/\theta}\).

For the subtracted tangent gradients, split the indices according as
\(A_i\leq1/2\) or \(A_i>1/2\).  In the first set
\(|d\log L_i|\leq C A_i|\II_i|\); in the second there are at most
\(4m\) indices.  Cauchy--Schwarz and
\eqref{eq:entropy-last-surplus} consequently bound the sum of their
norms by \(C_n\sqrt{mP_p/\theta}\).  The final tangent embedding sends
a vector \(v\) to \((v,(Ndf_i(v))_i)\).  Summing the recurrence, the
initial gradient has no vertical component in this common extension,
while the vertical component of the final gradient is
\((Ndf_i(b))_i\).
The normal and tangent bounds therefore prove
\[
 \sum_i(Ndf_i(b))^2\leq C_n(1+m)P_p/\theta.
\]
It follows that
\[
 |b(u)|\leq\frac C r|b(r)|
 \leq C_n\left(\frac{(1+m)P_p}{\theta N^2r}\right)^{1/2}
 \leq \frac{C_nP_p}{\sqrt{\theta E}}.
\]
For the last inequality use \(E\asymp\epsilon\), \(N^2=l\epsilon\),
and
\[
 P_p\geq\frac{c_n}{lr}
       \left(\epsilon+\frac{m^2}{\epsilon}\right)
       \geq \frac{c'_n(1+m)}{lr}
\]
for \(\epsilon\geq1\).  The same first surplus bound gives
\begin{equation}\label{eq:entropy-exit-derivatives}
 |du|^2\leq\frac{C_n}{N^2r}\leq\frac{C_nP_p}{E^2}.
\end{equation}
All these estimates are only needed where \(u\ne0\), a subset of
\(\{r\leq1/4\}\) for large \(E\).  Since
\(\theta E\asymp_n E/(1+\log E)^2\) tends to infinity, the two
derivative terms in \eqref{eq:entropy-conformal-exit} cost an
arbitrarily small fraction of \(P_p\).

There is enough surplus on this region also for the change in the
entropy term.  From \(r\leq E^{-0.32}\) and
\eqref{eq:entropy-last-surplus},
\[
 P\geq P_p\geq c_n E^{1.32}.
\]
Using \eqref{eq:entropy-before-exit} and \(e^{2u}\leq C E^{0.04}\),
we obtain there
\begin{equation}\label{eq:entropy-exit-volume-cost}
 \frac{e^{2u}W}{P}\leq
 C_nE^{0.04+0.26}P^{-0.25}\leq C_nE^{-0.03},
 \qquad
 \frac{e^{2u}nu}{P}\leq C_nE^{-1.28}\log E.
\end{equation}
In the first bound we used \(E+P\leq C_nP\), valid for large \(E\)
on this region.  Both quantities tend to zero with dimensional
constants.

We now assemble these estimates.  After replacing \(Y\) by \(Z\), write
the retained surplus as \(\alpha_n P\), where \(\alpha_n>0\) is fixed
dimensionally.  Put \(B_u=2|b(u)|+(n-1)|du|^2\).  Substituting in
\eqref{eq:entropy-conformal-exit} and moving the new entropy term
to the left gives
\begin{align*}
 &\wscal(\widehat h,\widehat\phi)
       -E e_{\widehat h}(Z)-2(W+nu)+n\log E\\
 &\quad\geq -(C+C_n)
   +e^{-2u}\bigl[\alpha_n P-B_u
            -2(e^{2u}-1)W-2nu e^{2u}\bigr]
   +(1-e^{-2u})n\log E .
\end{align*}
We used \(e^{-2u}(C+C_n)\leq C+C_n\).  The last term is nonnegative
for \(E\geq1\).  Equations \eqref{eq:entropy-exit-derivatives} and
\eqref{eq:entropy-exit-volume-cost} leave at least a fixed fraction
of \(P\) in the bracket on the support of \(u\); off that support the
volume and derivative costs vanish.  We may thus set
\[
 P'=\alpha'_n e^{-2u}P
\]
for a fixed \(\alpha'_n>0\), proving the scalar inequality in
\eqref{eq:entropy-after-exit}.  The only scalar constant charged was
the dimensional replacement cost.

The volume estimate at exit follows as well.  Since
\(P\leq C_ne^{2u}P'\leq C_nE^{0.04}P'\) and \(nu\leq C_n\log E\),
\[
 \widehat W\leq
 C_n E^{0.26+0.04(0.75)}(E+P')^{0.75}+C_n\log E
 \leq E^{0.32}(E+P')^{0.75}
\]
for all sufficiently large dimensional \(E\).

Finally, \eqref{eq:entropy-Z-energy} and the conformal scaling of
energy give
\[
 e_{\widehat h}(Z)^{1/2}
 \leq e^{-u}\frac{\sqrt n/N+c_EK_0}{\sqrt{r+c_E^2}}.
\]
The choice \eqref{eq:entropy-u-choice} implies
\[
 \frac{e^{-u}}{\sqrt{r+c_E^2}}
       \leq C\left(\frac{E^{-0.02}}{c_E}+E^{0.16}\right).
\]
Using \(\sqrt E/N\leq C_n(1+\log E)\), \(c_EK_0\leq E^{-1/2}\),
and \(N\leq\sqrt E\), we conclude that
\begin{equation}\label{eq:entropy-exit-gradient}
 e_{\widehat h}(Z)^{1/2}
 \leq C_n\left(
 \frac{(1+K_0)(1+\log E)}{E^{0.02}}+\frac{E^{0.16}}{N}\right).
\end{equation}
The first summand is \(o(E^{-0.01})\) for fixed \(K_0\); the second is
\(O_n((1+\log E)E^{-0.34})\).  Taking \(E\) large enough in terms of
\(n,K_0\) proves the final inequality in
\eqref{eq:entropy-after-exit}.  The functions \(u\) and the normalizing
denominator are smooth with finite bounds in each fixed construction.
The supports of \(Z\) are those of the last bumps and of \(Y^0\).
Thus the asserted uniformity and compact support properties persist.
\end{proof}

\subsection{A fixed graph and the final parameter choice}

\begin{proof}[Proof of Theorem~\ref{thm:entropy}]
Fix parameters \(H,A,B\) for the prescribed elementary bound and set
\(J=H+2\).  This number of phases is fixed before any phase target
is chosen.  The first target will be chosen after the bootstrap data
and their uniform constants have been fixed; the estimates below
hold for every sufficiently large choice of that target.

Apply Lemma~\ref{lem:entropy-bootstrap}.  At its output
\(W=\log(dV_{h_*}/\rho)=0\), \(\phi=0\), and
\(\epsilon_0\in[\epsilon_*,2\epsilon_*]\).  Increasing a dimensional
constant once, \eqref{eq:entropy-bootstrap-output} has the form
\[
 \wscal(h_*,0)\geq-C_n+2W-n\log\epsilon_0
                    +\epsilon_0 e_{h_*}(F_*),
\]
because \(\epsilon_0\) ranges in a fixed dimensional interval.
We use the schedule \eqref{eq:entropy-step-schedule} from now on.
The bound \eqref{eq:entropy-summability} applies to the entire
increasing sequence of energy coefficients, even when it is
partitioned into phases.

Choose a large first target and run the first phase until the parameter
first reaches it; call the actual output parameter \(E_1\).  The
overshoot is at most a factor \(1+l_*\).  In every sweep of this phase
use \(h_*\) as the reference metric.  The possible support enlargement
per sweep, including its separation buffers, has a dimensional radius
in this reference metric by Proposition~\ref{prop:sweep}.  The total
radius is therefore at most \(C_n(1+\log^3(1+E_1))\).

We record a bound for a single fixed graph sufficient for every
coloring in this phase and for all cooccurrences at its exit.  All
data before the choice of the first target, including \(h_*\) and its
uniform comparison with \(g\), have already been fixed.  The lifted
metric \(h_*\) is at least a positive input-dependent multiple of \(g\).
For a fixed finite generating set of \(\Gamma\), cocompactness implies
that word length is bounded above by an affine input-dependent
function of lifted \(g\)-distance between translates of fixed compact
fundamental sets.  For instance, subdividing a curve into bounded
length pieces and following the finitely many translates meeting
successive covering sets gives this bound.  The number of group
elements of word length at most \(L\) is at most exponential in \(L\).
It follows that the graph joining every pair of labels whose possible
enlarged supports in the first phase meet, or come within a required
buffer, has degree at most
\begin{equation}\label{eq:entropy-first-graph}
 \Delta_1+1\leq
 \exp\bigl(C_{\mathrm{inst}}(1+\log E_1)^3\bigr).
\end{equation}
Here \(C_{\mathrm{inst}}\) is finite and fixed \emph{before} the
first target is chosen.  It includes the initial number of label
orbits and the compact support bounds.  The two copies of labels
used at exit and their possible cooccurrences change this bound only
by a fixed factor, which is absorbed in \(C_{\mathrm{inst}}\).

Take \(E_1\) sufficiently large for
Lemma~\ref{lem:entropy-accumulation} and
Lemma~\ref{lem:entropy-phase-exit}; in the latter the initial \(K_0\)
is the finite input-dependent gradient bound of \(F_*\).
We now have the scalar inequality, the volume bound with exponent
\(0.32\), and the absolute gradient bound at \(E_1\).
Denote the exit metric, map, and fixed cooccurrence graph by
\(h_1,F_1,\mathcal G_1\).

We next isolate how a small gradient controls the graph in a later
phase.  Suppose a phase starts with metric \(h_j\), root map \(F_j\)
with \(e_{h_j}(F_j)^{1/2}\leq K\), and a fixed graph \(\mathcal G_j\)
containing all its possible cooccurrences.  Let \(R\) bound the
following pairwise distance: whenever two possible enlarged supports
meet or come close enough to require separation by a buffer, there
are points \(x,y\) in the old nonzero supports of their respective
parent labels with \(\dist_{h_j}(x,y)\leq R\).  If a closure is
involved, allow an arbitrarily small additional distance.  Suppose
\[
 KR<\sqrt2.
\]
Choosing the additional distance small enough, the gradient bound gives
\[
 \|F_j(x)-F_j(y)\|_{\ell^2}
 \leq K\dist_{h_j}(x,y)<\sqrt2.
\]
Two nonnegative unit vectors with disjoint supports have distance
exactly \(\sqrt2\).  Thus \(F_j(x)\) and \(F_j(y)\) share a nonzero
label, say \(k\).  The first parent and \(k\) cooccur at \(x\), and
the second parent and \(k\) cooccur at \(y\).  The two parents are
therefore at graph distance at most two in \(\mathcal G_j\).  This
argument applies separately at each transverse parameter; because
\(\mathcal G_j\) already contains all possible cooccurrences, the
conclusion holds in that one fixed graph.

After an exit at \(E_j\), we choose a target that both meets this
overlap condition and grows by one exponential level:
\begin{equation}\label{eq:entropy-phase-targets}
 T_{j+1}=\exp(E_j^{1/1000}),
\end{equation}
and call the first parameter reaching it \(E_{j+1}\).  Thus
\[
 T_{j+1}\leq E_{j+1}\leq(1+l_*)T_{j+1}.
\]
The choice \(1/1000\) leaves the strict inequality
\(3/1000<1/100\): the cubic logarithmic cost of support enlargement
will be smaller than the gain from the preceding gradient bound.
For large \(E_j\) this target satisfies
\(E_j^{0.32}\leq\tfrac12 E_{j+1}^{0.26}\).  The outgoing bound with
exponent \(0.32\) from phase \(j\) and
Lemma~\ref{lem:entropy-accumulation} therefore give the bound with
exponent \(0.26\) before the next exit.  The next initial gradient
bound is \(K_0\leq E_j^{-1/100}\leq1\), so all further thresholds
for Lemma~\ref{lem:entropy-phase-exit} are dimensional.  Increasing
the first target once ensures these conclusions in every one of a
prescribed finite number of phases.

Use \(h_j\), the starting metric, as the reference metric in phase
\(j+1\).  Let \(R_{j+1}\) bound the pairwise old-support distance in
the preceding argument, including both support enlargement radii and
all separation buffers.  The number of sweeps gives
\[
 R_{j+1}\leq C_n(1+\log^3 E_{j+1})
          \leq C_n E_j^{3/1000}.
\]
Consequently
\begin{equation}\label{eq:entropy-overlap-radius}
 R_{j+1}E_j^{-1/100}\leq C_n E_j^{-7/1000}\longrightarrow0.
\end{equation}
For large \(E_j\), the overlap condition therefore holds with
\(K=E_j^{-1/100}\).  For every relevant pair the preceding estimate
becomes
\[
 \|F_j(x)-F_j(y)\|_{\ell^2}
 \leq E_j^{-1/100}\dist_{h_j}(x,y)<\sqrt2.
\]
This uses only the starting metric of the phase, without a comparison
constant between that metric and \(g\).

It follows that every separation graph needed in phase \(j+1\) is
contained in the square \(\mathcal G_j^2\), after identifying labels
with their parents at the start of the phase.  At exit there are two
copies of each parent.  Joining two such copies whenever their parents
are equal or at distance at most two gives a fixed graph
\(\mathcal G_{j+1}\) containing all new cooccurrences.  If
\(\Delta_j\) is the degree of \(\mathcal G_j\), this construction gives
\begin{equation}\label{eq:entropy-graph-recurrence}
 \Delta_{j+1}+1\leq2(\Delta_j+1)^2.
\end{equation}
It also doubles the number of label orbits.  The preliminary coloring
construction applies to these graphs before each sweep.
All analytic support and chart bounds in the resulting finite
construction remain finite for the input, even though they need
not satisfy a useful quantitative bound.  They do not enter
\eqref{eq:entropy-graph-recurrence}.

After \(J\) phases, apply the coloring argument to
\(\mathcal G_J^2\), whose maximum degree is at most \(\Delta_J^2\).
Take \(q\) to be the maximum of its degree plus one, the number of
label orbits, and \(n+1\).  Equations \eqref{eq:entropy-first-graph} and
\eqref{eq:entropy-graph-recurrence} imply
\begin{equation}\label{eq:entropy-q-growth}
 \log q\leq C_{\mathrm{inst},J,n}(1+\log E_1)^3.
\end{equation}
For fixed \(J,n\) the coefficient here is fixed before \(E_1\) is
chosen.  The graph \(\mathcal G_J\) and this coloring have exactly
the properties asserted in the theorem.

It remains to make the first target large enough for the already fixed
\(J=H+2\).  Let
\(\mathcal B(q)\leq\exp^{[H]}(B(1+q)^A)\) be the prescribed
elementary bound.  Equation \eqref{eq:entropy-q-growth} implies
\[
 \mathcal B(q)\leq
 \exp^{[H+1]}\bigl(C'_{\mathrm{inst},J,n,A,B}
                         (1+\log E_1)^3\bigr)
\]
for large \(E_1\).  On the other hand, the recurrence
\eqref{eq:entropy-phase-targets} gives, for each fixed \(J\geq2\),
\begin{equation}\label{eq:entropy-tower-growth}
 E_J\geq\exp^{[J-1]}\bigl(c_J E_1^{1/1000}\bigr)
\end{equation}
for some \(c_J>0\) and all sufficiently large \(E_1\).
One proves this by induction: raising a fixed-height exponential
tower with argument \(c t\) to the power \(1/1000\) still dominates
a tower of the same height with argument \(c't\), for some \(c'>0\)
and all large \(t\); the next application of \(\exp\) adds a level.
For height one this follows from
\((e^{ct})^{1/1000}=e^{ct/1000}\); for higher heights, reducing the
bottom argument by any fixed factor absorbs the fixed outer
multiplicative factor.  This proves the induction.

With the already chosen \(J=H+2\), the tower heights on the
two sides are now the same, and
\[
 c_J E_1^{1/1000}>
 C'_{\mathrm{inst},J,n,A,B}(1+\log E_1)^3
\]
for all sufficiently large \(E_1\).  Thus \(E_J>\mathcal B(q)\).
We may choose the first target still larger to have
\(E_J>E_{\min}\), to satisfy the initial gradient threshold involving
\(K_0\), and to meet all dimensional thresholds used above.
These choices affect only the size of the final parameter and the
finite uniform bounds of the output objects.

Finally we account for the scalar constant.  The bootstrap charges
at most \(2C_b\epsilon_*\), independently of its input-dependent
number of sweeps.  The initial logarithmic normalization has a
dimensional cost because \(\epsilon_0\in[\epsilon_*,2\epsilon_*]\).
All entropy sweep losses together are bounded by
\(C_n\sum l^2\leq C_n\).  Each exit charges at most \(C_n\), and their
number \(J=H+2\) depends only on the prescribed elementary function
and \(n\).  The input gradient bound \(K_0\), the initial small
\(\epsilon_{\mathrm{in}}\), the compact support radii, and all metric
and elliptic bounds occur only in the choices of the first target,
graphs, and arbitrarily small analytic errors.  None contributes to
the scalar constant.  Discard the final nonnegative surplus \(P'\);
the remaining inequality is \eqref{eq:entropy-contract} with
\(C=C(n,\mathcal B)\).

The density \(\rho\) has been unchanged since the bootstrap, so
\(\rho\leq dV_g\).  The volume logarithm starts at zero and only
increases by \(W_p\geq0\) and \(nu\geq0\).  Finitely many sweeps and
exits preserve the asserted smoothness, uniformity, compact supports,
and finite number of free label orbits.  The transverse spaces were
enlarged only by the probability products used for the colorings,
and the sweep constructions are Borel on these spaces by
Proposition~\ref{prop:sweep}.  This proves all the conclusions.
\end{proof}
